%
% 3-hermite.tex
%
% (c) 2026 Prof Dr Andreas Müller
%
\section{Quadratfreie Faktorisierung und die Methode von Hermite
\label{buch:ratint:section:hermite}}
\kopfrechts{Quadratfreie Faktorisierung und die Methode von Hermite}%
Gemäss dem Plan am Ende von Abschnitt~\ref{buch:ratint:section:bernoulli}
ist, ist eine Faktorisierung des Nenners mit sehr speziellen
Teilbarkeitseigenschaften nötig.
Im vorangegangenen Abschnitt wurde mit dem euklidischen Algorithmus
das dazu geeignete Werzeug gefunden.
Dieses soll jetzt dazu eingesetzt werden, die Faktorisierung zu
finden und den Reduktionsschritt für die Integrale durchzuführen.

%
% Quadratfreie Faktorisierung
%
\subsection{Quadratfreie Faktorisierung}
Die quadratfreie Faktorisierung löst das Problem, den Nenner ohne
Bestimmung der Nullstellen so zu faktorisieren, dass eine
Partialbruchzerlegung möglich wird, mit der die Stammfunktion
bestimmt werden kann.

%
% Mehrfache Nullstellen
%
\subsubsection{Mehrfache Nullstellen}
Sei $q(z) \in \mathbb{C}[z]$ ein Polynom mit einer $n$-fachen Nullstelle
$\alpha$.
Dies bedeutet, dass $q(z)$ in der Form
\[
q(z)
=
(z-\alpha)^n\cdot s(z)
\]
mit einem Polynom $s(z)\in\mathbb{C}[z]$ geschrieben werden kann.
Das Polynom $s(z)$ hat $\alpha$ nicht als Nullstelle, $s(\alpha)\ne 0$,
es ist also nicht durch $z-\alpha$ teilbar.

Nach der Produktregel ist die Ableitung von $q(z)$
\[
q'(z)
=
n
(z-\alpha)^{n-1}
\cdot s(z)
+
(z-\alpha)^n\cdot s'(z)
=
(z-\alpha)^{n-1} 
\cdot (s(z)+(z-\alpha)\cdot s'(z)).
\]
Für $z=\alpha$ ist der Klammerfaktor
\[
s(\alpha) + (\alpha-\alpha)\cdot s'(\alpha)
=
s(\alpha)
\ne 0,
\]
der Klammerfaktor hat $\alpha$ nicht als Nullstelle.
Für eine mehrfache Nullstelle ist $n>1$ und es folgt, dass $q'(z)$
die Zahl $\alpha$ als $n-1$-fache Nullstelle hat.

Diese Beobachtung lässt sich auch verallgemeinern.
Hat $q(z)$ einen Faktor $p(z)$, der $n$-fach in $q(z)$ vorkommt,
dann kann man
\[
q(z)
=
p(z)^n\cdot s(z)
\]
schreiben, wobei $s(z)$ nicht durch $p(z)$ teilbar ist.
Die Ableitung von $q'(z)$ ist dann nach der Produktregel
\[
q'(z)
=
n p(z)^{n-1} \cdot s(z) + p(z)^n \cdot s'(z)
=
q(z)^{n-1}\cdot (ns(z) + p(z) s'(z)).
\]
Der Faktor auf der rechten Seite ist nicht durch $p(z)$ teilbar,
weil $s(z)$ nicht durch $p(z)$ teilbar ist.
Es folgt, dass $q'(z)$ durch $p(z)^{n-1}$ teilbar ist.

Wenn also $q(z)$ einen mehrfachen vorkommenden Faktor hat,
dann ist dieser Faktor ein gemeinsamer Teiler von $q(z)$
und $q'(z)$.

\begin{beispiel}
TODO
\end{beispiel}

%
% Quadratfreie Polynome
%
\subsubsection{Quadratfreie Polynome}
Damit $q(z)$ und $q'(z)$ keinen nichtrivialen gemeinsamen Teiler
haben, darf jeder Linearfaktor nur ein einziges Mal vorkommen.
Solche Polynome heissen quadratfrei im Sinne der folgenden
Definition.

\begin{definition}
Ein Polynom in $\mathbb{C}[z]$ heisst \emph{quadratfrei}, wenn es
keine mehrfachen Nullstellen hat oder wenn es keine mehrfachen
Faktoren hat.
\index{quadratfrei}%
\end{definition}

Über den komplexen Zahlen zerfällt also ein quadratfreies Polynom
immer in ein Produkt von \emph{verschiedenen} Linearfaktoren.

Wenn ein Polynome $q(z)$ mehrfache Nullstellen hat, dann hat die
Ableitung die gleichen Nullstellen, aber mit einer um $1$ kleineren
Ordnung.
Der grösste gemeinsame Teiler von $q(x()$ und $q'(x)$ besteht
daher alle Faktoren, die mehr als einmal vorkommen, mit einer
gegenüber $q(z)$ um $1$ verminderten Vielfachheit.
Im Quotienten
\[
p(z)
=
\frac{q(z)}{\operatorname{ggT}(q(z),q'(z))}
\]
bleiben dann nur noch alle Linearfaktoren genau einmal stehen.

\begin{beispiel}
Das Polynom 
\[
q(z)
=
z^5-z^4-2z^3+2z^2+z-1
\]
hat die Ableitung
\[
q'(z)
=
5z^4-4z^3-6z^2+4z+1.
\]
Der euklidische Algoritmus liefert als grössten gemeinsamen
Teiler das Polnom
\[
g(z)
=
\operatorname{ggT}(q(z), q'(z))
=
z^3-z^2-1+1.
\]
Der Quotient
\[
\frac{q(z)}{g(z)}
=
z^2-1
\]
ist ein quadratfreies Polynom, welches alle Linearfaktoren von $q(z)$
enthält.
Es ist in diesem Fall einfach zu sehen, dass $z^2-1=(z-1)(z+1)$,
d.~h.~die einzigen Nullstellen von $q(x)$ sind $\pm 1$.

Man kann aber die Faktoren auch noch auf etwas systematischere Art
finden.
Dazu berechnet man den grössten gemeinsamen Teiler von $g(z)$ und $g'(z)$,
wegen $g'(z) = 3z^2 -2z -1$ ist
ist
\[
\operatorname{ggT}(g(z),g'(z))
=
z-1
\]
und
\[
\frac{g(z)}{\operatorname{ggT}(g(z),g'(z))}
=
z^2-1.
\]
Daraus kann man schliessen, dass es zwei Linearfaktoren von $q(z)$
gibt.
Einer davon kommt genau zweimal vor, der andere ist $z-1$ und kommt
dreimal vor.
Durch Division $(z^2-1)/(z-1)=z+1$ findet man jetzt, dass $z+1$
der Faktor ist, der genau zweimal ist.
Wir haben somit die Faktorisierung
\[
q(z)
=
(z+1)^2 (z-1)^3
\]
gefunden.
\end{beispiel}

Das Beispiel zeigt, dass die Suche nach quadratfreien Faktoren
mithilfe des grössten gemeinsamen Teilers manchmal sogar eine
Faktorisierung des Polynoms liefern kann.

%
% Die Aufgabe der Faktorisierung in quadratfreie Faktoren
%
\subsubsection{Die Aufgabe der Faktorisierung in quadratfreie Faktoren}
Ein beliebiges normiertes Polynom $q(z)$ kann einzelne Nullstellen
mehrfach enthalten.
Man kann aber immer noch eine Faktorisierung der Form
\[
q(z)
=
\prod_{k=1}^n
(z-\alpha_k)^{n_k},
\]
wobei alle Nullstellen $\alpha_k$ verschieden sind und $n_k$ ihre
Vielfachheit ist.
Fasst man die Faktoren mit gleicher Vielfachheit zusammen, kann man
zu jeder Vielfachheit $l$ das Polynom
\[
q_l(z) = \prod_{n_k=l} (z-\alpha_k)
\]
bilden.
Da die Faktoren von $q_l(z)$ im Polynom $q(z)$ genau $l$-mal vorkommen,
kann man das Polynom $q(z)$ aus
\[
q(z)
=
q_1(z)
\cdot
q_2(z)^2
\cdot
q_3(z)^3
\cdot
\ldots
\cdot
q_m(z)^m
\]
rekonstruieren.
Nach Konstruktion kann ein Linearfaktor nur immer in genau einem
der Faktoren $q_l(z)$ vorkommen, die Faktoren sind also paarweise
teilerfremd.

\begin{definition}[quadratfreie Faktorisierung]
\index{quadratfreie Faktorisierung}%
\index{Faktorisierung, quadratfrei}%
Die \emph{quadratfreie Faktorisierung} eines Polynoms $q(z)$ ist eine
Faktorisierung der Form
\[
q(z)
=
a
\prod_{l=1}^m q_l(z)^l
\]
mit teilerfremden Faktoren.
\end{definition}

Die Bestimmung der quadratfreien Faktorisierung aus der Zerlegung
in Linearfaktoren zeigt zwar, dass es eine solche Faktorisierung
gibt, ist aber kein praktikabler Weg, sie algorithmisch zu bestimmen.
daher soll in diesem Abschnitt soll die folgende Aufgabe gelöst werden.

\begin{aufgabe}
Formuliere einen Algorithmus, der zu einem beliebigen rationalen Polynom
$q(z)\in\mathbb{Q}[z]$ die quadratfreie Faktorisierung findet.
\end{aufgabe}

%
% Berechnung der quadratfreien Faktorisierung
%
\subsubsection{Berechnung der quadratfreien Faktorisierung}
Aus dem Beispiel lesen wir ab, dass der grösste gemeinsame
Teiler $g(z)=\operatorname{ggT}\bigl(q(z),q'(z)\bigr)$ von $q(z)$
und $q'(z)$ alle mehrfachen Faktoren enthält und der Quotient
$q(z)/g(z)$ alle Linearfaktoren genau einmal.
Diese Eigenschaft soll jetzt dazu verwendet werden, die Faktoren
$q_l(z)$ zu finden.
Dazu konstruieren wir zunächst zwei Folgen $g_k(z)$ und $p_k(z)$
wie folgt.
Das Polynom $g_0(z)$ ist $g_0(z)=q(z)$.
Die Polynome $g_k(z)$ und $p_k(z)$ mit $k>0$ sind rekursiv definiert
durch
\begin{align*}
g_{k+1}(z) &= \operatorname{ggT}\bigl(g_k(z),g'_k(z)\bigr)
&&\text{und}&
p_{k+1}(z) &= \frac{g_k(z)}{g_{k+1}(z)}.
\end{align*}
Nach Konstruktion sind die Polynom $p_k(z)$ alle quadratfrei.
Sie enthalten alle Linearfaktoren von $g_k(z)$.
Nach Konstruktion ist $g_{k+1}(z)$ immer ein echter Teiler von $g_k(z)$.
Der Grad nimmt daher in jedem Schritt ab.
Die Konstruktion bricht daher nach endlich vielen Schritten ab.

In jedem Iterationsschritt nimmt der Exponent der in $g_k(z)$
vertretenen Linearfaktoren um $1$ ab.
Das Ende $g_k(z)=1$ wird also genau in dem Moment erreicht, 
in dem der grösste Exponent auf $1$ reduziert worden ist.
Somit ist $k$ der grösste Exponent eines Linearfaktors.

Die Polynome $p_k(z)$ enthalten alle Linearfaktoren, die
mindestens $k$ mal in $q(z)$ aufreten.
Es folgt, dass $p_{k+1}(z)$ ein Teiler von $p_k(z)$ ist und dass
der Quotient $p_k(z)/p_{k+1}(z)$ aus den Faktoren besteht, die 
genau $k$-mal in $q(z)$ vorkommen.
Dies ist genau, was die Polynome $q_l(z)$ definiert.
Dies beweist den folgenden Satz.

\begin{satz}[quadratfreie Faktorisierung]
Sei $q(z)\in K[z]$ ein normiertes Polynom.
Setze $q_0(z)=q(z)$ und bilde rekursiv die Folgen
\begin{align*}
g_k(z) &= \operatorname{ggT}\bigl(g_{k-1}(z),g_{k-1}'(z)\bigr)
&
p_k(z) &= \frac{g_{k-1}(z)}{g_k(z)}
&
q_k(z) &= \frac{p_k(z)}{p_{k+1}(z)}.
\end{align*}
Sei $m$ dasjenige $k$, für das $p_k(z)=1$ ist.
Dann ist
\[
q(z)
=
\prod_{l=1}^m q_l(z)^l
\]
die quadratfreie Faktorisierung.
\end{satz}

\begin{beispiel}
Wir betrachten das Polynom
\begin{align*}
q(z)
&=
z^9+z^8-5z^7-5z^6+9z^5+9z^4-7z^3-7z^2+2z+2
\intertext{mit der Ableitung}
q'(z)
&=
9z^8+8z^7-35z^6-30z^5+45z^4+36z^3-21z^2-14z+2
\end{align*}
Die Polynome $g_k(z)$, $p_k(z)$ und $q_l(z)$ sind
\begin{align*}
g_1(z)
&=
z^5+z^4-2z^3-2z^2+z+1,
&
p_1(z)
&=
z^4-3z^2+2,
&
q_1(z)
&=
z^2-2,
\\
g_2(z)
&=
z^3+z^2-z-1,
&
p_2(z)
&=
z^2-1,
&
q_2(z)
&=
1,
\\
g_3(z)
&=
z+1
&
p_3(z)
&=
z^2-1,
&
q_3(z)
&=
z-1,
\\
g_4(z)
&=
1,
&
p_4(z)
&=
z+1,
&
q_4(z)
&=
z+1,
\\
g_5(z)
&=
1,
&
p_5(z)
&=
1.
&
&
\end{align*}
Tatsächlich kann man nachrechnen, dass
\[
q(z)
=
(z^2-2)(z-1)^3(z+1)^4
\]
ist.
Man beachte aber, dass diese Faktorisierung mit $z^2-2$ auch einen Faktor
enthält, der kein Linearfaktor ist.
Da $z^2-2 = (z+\!\sqrt{2})(z-\!\sqrt{2})$, gibt es über $\mathbb{Q}$
gar keine Faktorisierung in Linearfaktoren.
\end{beispiel}

\begin{beispiel}
\label{buch:ratint:hermite:bsp:qf2}
Man finde die quadratfreie Faktorisierung des Polynoms
\[
q(z) = 9z^6 + 6z^5 - 65z^4 + 20z^3 + 135z^2 - 154z + 49
\]
mit der Ableitung
\[
q'(z) = 54z^5+30z^4-260z^3+60z^2+270z-154.
\]
Die Polynome $g_k(z)$, $p_k(z)$ und $q_l(z)$ sind
\begin{align*}
g_1(z)
&= 
3z^4-2z^3-12z^2+18z-7,
&
p_1(z)
&=
3z^2+4z-7,
&
q_1(z)
&=
1,
\\
g_2(z)
&=
z^2-2z+1,
&
p_2(z)
&=
3z^2+4z-7,
&
q_2(z)
&=
3z+7,
\\
g_3(z)
&=
z-1,
&
p_3(z)
&=
z-1,
&
q_3(z)
&=
1,
\\
g_4(z)
&=
1,
&
p_4(z)
&=
z-1,
&
q_4(z)
&=
z-1,
\\
g_5(z)
&=
1,
&
p_5(z)
&=
1.
&&
\end{align*}
Daraus liest man die quadratfreie Faktorisierung
\[
q(z)
=
(3z+7)^2(z-1)^4
\]
ab.
\end{beispiel}

%
% Reduktion des Integrals
%
\subsection{Reduktion des Integrals}
Mit dem Werkzeug der quadratfreien Faktorisierung lässen sich jetzt
die für den rationalen Teil des Integrals nötigen Schritte des Plans
vom Ende von Abschnitt~\ref{buch:ratint:section:bernoulli} durchführen.
Zu diesem Zwck muss erst eine Partialbruchzerlegung gefunden werden.
Anschliessend müssen die einzelnen Terme der Partialbruchzerlegung
integriert werden.

%
% Partialbruchzerlegung zur quadratfreien Faktorisierung
%
\subsubsection{Partialbruchzerlegung zur quadratfreien Faktorisierung}
Der wesentliche Schritt zur Konstruktion der Partialbruchzerlegung
ist der folgende Satz.

\begin{satz}
Sei 
\[
f(z)
=
\frac{p(z)}{q(z)^m r(z)},
\]
wobei $q(z)$ teilerfremd zu $r(z)$ ist.
Dann gibt es 
\[
f(z)
=
f_0(z)
+
\frac{s(z)}{r(z)}
+
\sum_{k=1}^m \frac{t_k(z)}{q(z)^k},
\]
wobei $f_0(z)$ ein Polynom ist, $s(z)$ ein Polynom mit Grad
$\deg s(z)<\deg r(z)$ und die Polynome $t_k(z)$ Grad $\deg t_k(z)<\deg q(z)$
haben.
\end{satz}

\begin{proof}
Da $q(z)$ und $r(z)$ teilerfremd sind, sind auch $q(z)^m$ und $r(z)$
teilerfremd.
Daher gibt es $s(z)$ und $t(z)$ mit
\begin{equation}
p(z)
=
s(z)q(z)^m + t(z)r(z).
\label{buch:ratint:hermite:eqn:p=sq+tr}
\end{equation}
Einsetzen in den Bruch ergibt
\begin{align}
f(z)
&=
\frac{p(z)}{q(z)^m r(z)}
=
\frac{s(z)q(z)^m + t(z)r(z)}{q(z)^mr(z)}
=
\frac{s(z)}{r(z)}
+
\frac{t(z)}{q(z)^m}.
\label{buch:ratint:hermite:eqn:srtq}
\end{align}
Falls der Zähler einen Grad mindestens so gross wie der Nenner
hat, kann durch Division des Zählers mit Rest das Polynom $f_0(z)$
abgespaltet werden.

Es muss jetzt nur noch gezeigt werden, dass das der zweite Bruch
\eqref{buch:ratint:hermite:eqn:srtq}
in eine Summe der verlangten Art umgefomrt werden kann.
Dazu schreibt man den Zähler als
\begin{equation*}
t(z)
=
a_m(z) q(z) + t_{m}(z)
\end{equation*}
mit $\deg b(z) < \deg q(z)$.
Dann wird
\begin{align*}
\frac{t(z)}{q(z)^m}
&=
\frac{a(z)q(z)+b(z)}{q(z)^m}
=
\frac{a(z)}{q(z)^{m-1}}
+
\frac{b(z)}{q(z)^m}.
\intertext{Durch wiederholte Anwendung dieser Beobachtung auf
den neuen Zähler $a_m(z) = a_{m-1}(z)q(z)+t_{m-1}(z)$
wird daraus eine Summe der Form}
&=
\frac{t_1(z)}{q(z)}
+
\frac{t_2(z)}{q(z)^2}
+
\ldots
+
\frac{t_m(z)}{q(z)^m}.
\end{align*}
Damit ist der Satz vollständig bewiesen.
\end{proof}

Durch Anwendung auf eine Faktorisierung von $r(z)$ kann schrittweise
die Partialbruchzerlegung mit den Teilnennern $q_l(z)^k$, $k\le l$
für alle $l$ gefunden werden.
Man beachte, dass dazu wieder nur die Polynomdivision
und für die Identität 
\eqref{buch:ratint:hermite:eqn:p=sq+tr}
der euklidische Algorithmus benötigt wird.
Beide können für Polynome in einer beliebigen Körpererweiterung
für den Koeffizientenkörper durchgeführt werden.

\begin{beispiel}
\label{buch:ratint:hermite:bsp:partialbruch}
Man finde die Partialbruchzerlegung von
\[
\frac{r(z)}{q(z)}
\quad\text{mit}\quad
\left\{\;
\renewcommand{\arraycolsep}{2pt}
\begin{array}{rcl}
r(z) &=& 6651z^4 + 3969z^2 - 37338z + 21854 \\
q(z) &=& 9z^6 + 6z^5 - 65z^4 + 20z^3 + 135z^2 - 154z + 49.
\end{array}\right.
\]
Man verwendet dafür die quadratfreie Faktorisierung
\[
q(z)
=
(z-1)^4(3z+7)^2,
\]
die in Beispiel~\ref{buch:ratint:hermite:bsp:qf2} gefunden worden war.
Zunächst sind die beiden Faktoren $(z-1)^4$ und $(3z+7)^2$
teilerfremd, so dass es Polynome $s(z)$ und $t(z)$ gibt derart,
dass
\begin{align*}
r(z) &= s(z)(z-1)^4 + t(z) (3z+7)^2
\intertext{nämlich}
s(z) &=
	\frac{1}{50000}(3214890^5 z + 10180485 z^4 + 1928934 z^3 - 12037977 z^2
	- 46842138 z\\
     &\phantom{50000}\qquad + 33633306)
\\
t(z) &=
	\frac{1}{50000}(357210 z^7 - 1964655 z^6 + 5056506 z^5 - 9737280 z^4
	+ 15174810 z^3\\
     &\phantom{50000}\qquad - 37747395 z^2  + 52924410 z - 21613606),
\end{align*}
wie man mit dem erweiterten euklidischen Algorithmus finden kann.
Bildet man jetzt die Quotienten $s(z)/q(z)$ und $t(z)/q(z)$, heben sich die
polynomiellen Teile weg und es bleibt
\[
\frac{r(z)}{q(z)}
=
\frac{2646}{(3z+7)^2}
+
\frac{441z^2-882z+392}{(z-1)^4}.
\]
Der zweite Term hat noch nicht die von der Partialbruchzerlegung verlangte
Form.
Dies wird erreicht, indem man den Zähler wiederholt mit Rest durch $z-1$
dividiert:
\[
\left.
\begin{aligned}
441z^2-882z+392 &= (441z-441) (z-1 ) - 49 \\
441z-441 &= 441 (z-1) + 0
\end{aligned}\right\}
\Rightarrow
441z^2-882z+392 = 441 (z-1)^2 - 49.
\]
Daraus ergeben sich jetzt die Partialbrüche zum Teilnenner $z-1$ als
\[
\frac{t(z)}{(z-1)^4}
=
\frac{441}{(z-1)^2} - \frac{49}{(z-1)^4}.
\]
Insgesamt ist also
\[
\frac{r(z)}{q(z)}
=
\frac{2646}{(3z+7)^2}
+
\frac{441}{(z-1)^2} - \frac{49}{(z-1)^4}
\]
die gesuchte Partialbruchzerlegung.
\end{beispiel}


%
% Integration des rationalen Teils
%
\subsubsection{Integration des rationalen Teils}
Ganz analog der Vorgehensweise in der Methode von Bernoulli im
Abschnitt~\ref{buch:ratint:section:bernoulli} können auch die
Integranden mit Potenzen quadratfreier Polynome $q(z)$ im Nenner
auf den Nenner $q(z)$ reduziert werden.

\begin{satz}[Hermite]
\label{buch:ratint:hermite:satz:reduktion}
Sei $q(z)$ ein quadratfreies Polynom und $p(z)$ ein beliebiges
Polynom von kleinerem Grad $\deg q(z)<\deg p(z)$.
Dann gibt es Polynome $s(z)$ und $t(z)$ Polynome mit
\begin{equation}
s(z)\cdot q(z)
+
t(z)\cdot q'(z)
=
p(z),
\label{buch:ratint:hermite:ggtsz}
\end{equation}
für die zudem
\begin{align}
\int
\frac{p(z)}{q(z)^m}
\,dz
&=
\int
\frac{s(z)+t'(z)/(m-1)}{q(z)^{m-1}}
\,dz
+
\frac{t(z)}{(1-m)q(z)^{m-1}}
\label{buch:rating:hermite:satz:hermite:eqn:int}
\end{align}
wird.
\end{satz}

\begin{proof}
Da $q(z)$ quadratfrei ist sind  $q(z)$ und $q'(z)$ teilerfremd.
Nach dem euklidischen Algorithmus für Polynome gibt es daher
Polynome $s(z)$ und $t(z)$ derart, dass \eqref{buch:ratint:hermite:ggtsz}
gilt.
Einsetzen ins Integral ergibt
\begin{align*}
\int
\frac{p(z)}{q(z)^m}
\,dz
&=
\int
\frac{s(z)\cdot q(z)+t(z)\cdot q'(z)}{q(z)^m}
\,dz
\\
&=
\int
\frac{s(z)}{q(z)^{m-1}}
\,dz
+
\int
t(z)
\cdot
\frac{q'(z)}{q(z)^m}
\,dz
\intertext{Auf das zweite Integral kann partielle Integration
angewendet werden.
Das Polynom $p(z)$  wird abgeleitet während der Bruch
$q'(z)/q(z)^m$ integriert wird.
So entsteht}
&=
\int
\frac{s(z)}{q(z)^{m-1}}
\,dz
+
t(z)
\frac{1}{(1-m)q(z)^{m-1}}
-
\int
\frac{t'(z)}{(1-m)q(z)^{m-1}}
\,dz
\\
&=
\int
\frac{s(z)+t'(z)/(m-1)}{q(z)^{m-1}}
\,dz
+
t(z)
\frac{1}{(1-m)q(z)^{m-1}}.
\end{align*}
Damit ist der Satz bewiesen.
\end{proof}

Durch wiederholte Anwendung des
Satzes~\ref{buch:ratint:hermite:satz:reduktion}
kann man die Potenz des Nenners immer weiter reduzieren, bis
nur noch das Integral
\[
\int
\frac{p(z)}{q(z)}
\,dz
\]
bestimmt werden muss, wobei $p(z)$ und $q(z)$ teilerfremd sind,
$\deg p(z)<\deg q(z)$ und $q(z)$ quadratfrei.
Dieser Teil des Integrationsproblems heisst auch der 
\emph{logarithmische Teil} des Integrals.
\index{logarithmischer Teil}%
Die Methode von Hermite vermeidet eine vollständige Faktorisierung
und ersetzt sie durch eine quadratfreie Faktorisierung.
Sie erreicht aber auch nur eine Reduktion auf quadratfreie Nenner.
Für die vollständig Berechnung der Stammfunktion ist nach der Methode
von Bernoulli weiterhin eine Faktorisierung der quadratfreien Faktoren
sein.
Dies wird aber potentiell dadurch vereinfacht, dass alle Nullstellen
eines quadratfreien Polynoms einfach sind.

\begin{beispiel}
Man finde die Stammfunktion
\begin{align*}
F(z)
&=
\int
\frac{p(z)}{q(z)}\,dz
\\
&=
\int
\frac{
441z^7 + 780z^6 - 2861z^5 + 4085z^4 + 7695z^3 + 3713z^2 - 43253z + 24500
}{
9z^6 + 6z^5 - 65z^4 + 20z^3 + 135z^2 - 154z + 49
}
\,dz
\end{align*}
mit den Polynomen
\begin{align*}
p(z)
&=
441z^7 + 780z^6 - 2861z^5 + 4085z^4 + 7695z^3 + 3713z^2 - 43253z + 24500
\\
q(z)
&=
9z^6 + 6z^5 - 65z^4 + 20z^3 + 135z^2 - 154z + 49.
\end{align*}
Zunächst muss man den Bruch mit Reset dividieren und erhält
\[
F(z)
=
\int 49z + 54\,dz
+
\int
\frac{6651z^4 + 3969z^2 - 37338z + 21854}{q(z)}\,dz,
\]
wobei wir den Zähler im zweiten Integral mit $r(z)$ bezeichnen wollen.
Für die weitere Vereinfachung des zweiten Integrals wird jetzt die
quadratfreie Faktorisierung von $q(z)$ benötigt.
Sie ist
\[
q(z)
=
(z-1)^4(3z+7)^2.
\]
In diesem Fall hat die quadratfreie Faktorisierung eine vollständige
Faktorisierung in Linearfaktoren geliefert.
Damit kann jetzt auch die Partialbruchzerlegung gefunden werden,
wie dies in Beispiel~\ref{buch:ratint:hermite:bsp:partialbruch}
bereits geschehen ist.
Sie ist
\begin{equation}
\frac{r(z)}{q(z)}
=
\frac{294}{(z+\frac{7}{3})^2}
+
\frac{441}{(z-1)^2}
-
\frac{49}{(z-1)^4}.
\label{buch:ratint:hermite:bps:hermite:pf}
\end{equation}
Für die einzelnen Summanden ist jetzt noch die Reduktion nach der Methode
von Hermite von Satz~\ref{buch:ratint:hermite:satz:reduktion} durchzuführen.

Für den ersten Term müssen mit dem euklidischen Algorithmus
$s(z)$ und $t(z)$ gefunden werden derart, dass
\[
s(z)(z+{\textstyle\frac{3}{7}}) + t(z)(z+{\textstyle\frac{3}{7}})
=
s(z)(z+{\textstyle\frac{3}{7}}) + t(z) = 294
\]
gilt, nämlich $s(z)=0$ und $t(z)=294$.
Daraus folgt das Integral für den ersten Term als
\[
\int \frac{294}{(z+\frac37)^2}
=
-\frac{294}{z+{\textstyle\frac37}},
\]
wie man auch von der bekannten Integrationsformel erwartet hätte.

Für den zweiten Term muss man $s(z)$ und $t(z)$ finden derart, dass
\[
s(z)(z-1) + t(z)(z-1)' = s(z)(z-1) + t(z) = 441
\]
gilt.
Der euklidische Algorithmus liefert $s(z)=0$ und $t(z)=441$ und daher
\[
\int \frac{441}{(z-1)^2}\,dz
=
-
\frac{441}{z-1}.
\]
Ebenso funktioniert die Methode für den dritten Term.
Man muss $s(z)$ und $t(z)$ finden derart, dass
\[
s(z)(z-1) + t(z)(z-1)' = s(z)(z-1) + t(z) = -49
\]
gilt.
Es folgt $s(z)=0$ und $t(z)=-49$ und damit 
\[
\int \frac{-49}{(z-1)^4} \,dz
=
\frac{49}{3(z-1)^3}.
\]
Damit ist das gesuchte Integral
\[
F(z)
=
-
\frac{294}{z+\frac{3}{7}}
-
\frac{441}{z-1}
+
\frac{49}{3(z-1)^3}
\]
gefunden.
\end{beispiel}

\begin{beispiel}
Man bestimme das Integral
\[
F(z)
=
\int \frac{p(z)}{q(z)} \,dz
=
\int \frac{z+1}{(z^2+1)^3} \,dz.
\]
Der Nenner liegt bereits in der quadratfreien Faktorisierung mit dem
Faktor $q(z)=z^2+1$ vor, somit hat der Integrand bereits die Form der
Partialbruchzerlegung, es bleibt nur noch, die Reduktion nach Hermite 
mit $m=3$.
Dazu werden zunächst $s(z)$ und $t(z)$ bestimmt mit
\[
s(z)q(z) + t(z)q'(z) = z+1,
\]
man findet $s(z)=z+1$ und $t(z)=-\frac12(z^2+z)$.
Einsetzen in die Formel~\eqref{buch:rating:hermite:satz:hermite:eqn:int}
ergibt
\begin{align}
\int
\frac{z+1}{(z+1)^3}
\,dz
&=
\int \frac{2z}{3(z^2+1)^2} \,dz
+
\frac{z^2+z}{4z^4+8z^2+4}.
\label{buch:ratint:hermite:bsp:hermite2:eqn:i3}
\end{align}
Der Prozess muss jetzt für das verbleibende Integral auf der rechten
Seite wiederholt werden.
Diesmal findet man
\[
s(z)=\frac{2z+3}{4}
\qquad\text{ und }\qquad
t(z)=-\frac{2z^2+3}{8}
\]
und nach Einsetzen
\begin{align}
\int \frac{2z}{3(z^2+1)^2} \,dz
&=
\int \frac{\frac{3}{8}}{z^2+1}\,dz
+
\frac{2z^2+3z}{8z^2+8}.
\label{buch:ratint:hermite:bsp:hermite2:eqn:i2}
\end{align}
Kombiniert man
\eqref{buch:ratint:hermite:bsp:hermite2:eqn:i3}
und
\eqref{buch:ratint:hermite:bsp:hermite2:eqn:i2}
findet man
\begin{align*}
\int
\frac{z+1}{(z+1)^3}\,dz
&=
\frac38 \int \frac{1}{z^2+1}\,dz
+
\frac{2z^4+3z^3+4z^2+5z}{8z^4+16z^2+8}
\\
&=
\arctan z
+ 
\frac{3z^3+5z-2}{8z^4+16z^2+8}
+\frac14.
\end{align*}
Der Term $\frac14$ auf der echten Seite ist entstanden durch Division
des Zählers mit Rest durch den Nenner.
Er kann in der Integrationskonstanten absorbiert werden.
\end{beispiel}

%
% Die Methode von Horowitz-Ostrogradski
%
\subsection{Die Methode von Horowitz-Ostrogradski}
Die Berechnung der quadratfreien Zerlegung sowie die Reduktion
der Nennerpotenz mit der Methode von Hermite ist zwar algebraisch
einfach, verlangt aber die häufig wiederholte Berechnung von 
grössten gemeinsamen Teilern von Polynomen mit dem erweiterten
euklidischen Algorithmus.
Da jetzt aber die erwartete Form des rationalen und des logarithmischen
kennen, können die zugehörigen Zählerpolynome auch direkt mithilfe
eines linearen Gleichungssystems gesucht werden.
Dies ist die Methode von Horowitz, die allerdings schon früher von
Ostrogradski gefunden wurde.


